\chapter{Slips, Trips and Falls Risk Assessment}
\label{chap:Slips, Trips and Falls Risk Assessment}

%---------------------------- SECTION ----------------------------
\section{Why this assessment matters in construction}

Slips, trips and falls on the same level are among the most common causes of non-fatal injury in the UK construction industry, accounting for over a third of all major injuries reported to the HSE annually. The word ``common'' tends to diminish the seriousness of this hazard category: same-level falls on construction sites produce broken wrists, knee ligament ruptures, fractured ankles, head injuries from contact with protruding steel or concrete, and lacerations from falls onto sharp debris. The costs in human welfare, sickness absence, and civil liability are substantial, and the pattern of incidents is stubbornly resistant to improvement because the hazards are treated as inevitable features of the construction environment rather than manageable risks.

The construction site environment is intrinsically more challenging for pedestrian movement than almost any other workplace. Ground surfaces are routinely contaminated with mud, water, cement slurry, oil, and construction debris. Temporary access routes are poorly planned or absent. Excavations and level changes are unmarked. Cables, hoses, and pipe runs cross pedestrian routes without management. Materials are stored in circulation routes. Lighting is inadequate in covered or enclosed areas. These conditions are not inherent features of construction work; they are the products of inadequate site organisation, inadequate planning of pedestrian routes, and inadequate housekeeping, all of which can be addressed through the risk assessment process.

The particular vulnerability of the construction site environment to slip, trip and fall incidents arises from the interaction between changing ground conditions and the concentration of attention that operatives direct towards their work tasks. A bricklayer ascending to a scaffold platform with a hod of bricks is not watching the ground. A groundworker stepping backwards while laying pipe is not watching where they are stepping. The assessment must therefore identify hazards in the paths that operatives will actually take in the course of their work, not just the designated pedestrian routes, and must implement controls that do not depend on operatives continuously attending to the ground beneath their feet.

Footwear selection is frequently cited in post-incident investigations as a contributing factor: flat-soled safety boots worn on wet clay or smooth concrete, trainers worn by sub-contractors who have not been inducted to footwear standards, and unsuitable footwear worn during delivery operations are all recurring findings. The risk assessment must address footwear requirements as part of the control framework rather than leaving them to individual operative choice.

\barquote{``A tidy site is a safe site --- not because tidiness is an aesthetic virtue but because it removes the hazards that cause the injuries that end careers.''}

\example{Site Reality}{A joinery subcontractor is installing internal doors on the third floor of a new hotel development. The general contractor has completed plastering and screed laying but has not cleared the plastic screed mesh residues, wire ties, and off-cut timber from the circulation corridor. The joinery operative, carrying a door set, catches his boot on a protruding length of screed mesh, stumbles, and cannot arrest his fall because both hands are supporting the door. He falls forward, sustaining a fractured wrist and laceration to his forehead from contact with a protruding wall anchor bolt. The joinery contractor's RAMS identified ``slips and trips from construction debris'' as a hazard; the control specified was ``good housekeeping''. No one had defined whose responsibility the housekeeping was, what standard was required, or who was responsible for inspecting the corridor before work commenced.}

%---------------------------- SECTION ----------------------------
\section{Legal and professional framework}

The principal legislative duty for slips, trips and falls in construction is derived from the Health and Safety at Work etc.\ Act 1974, the Management of Health and Safety at Work Regulations 1999 (suitable and sufficient risk assessment), and the Construction (Design and Management) Regulations 2015 (which require the principal contractor to manage and coordinate welfare and housekeeping on the common parts of the site). The Workplace (Health, Safety and Welfare) Regulations 1992 apply to construction sites used as workplaces and require that floors and traffic routes are suitable, free from obstructions, and adequately lit; while these Regulations are primarily directed at fixed workplaces, the HSE applies their principles to construction site traffic routes.

HSE publication HSG150 (Health and Safety in Construction) provides sector-specific guidance on site organisation, housekeeping, and pedestrian access management. INDG293 (Preventing Slips and Trips at Work) provides the general framework for slip and trip risk management. CDM 2015 requires the Construction Phase Plan to address welfare provisions and site organisation, including the management of pedestrian and vehicle interfaces and the condition of access routes.

\paragraph{Key frameworks relevant to this chapter include: MHSWR 1999; CDM 2015; Workplace (Health, Safety and Welfare) Regulations 1992; HSG150; INDG293.}

\begin{itemize}
  \bitem{CDM 2015}{Requires the principal contractor to plan, manage, monitor, and coordinate the construction phase so that welfare and access arrangements protect the health and safety of all persons on the construction site; site organisation including pedestrian route management, housekeeping standards, and lighting of access routes must be addressed in the Construction Phase Plan.}
  \bitem{Workplace (Health, Safety and Welfare) Regulations 1992}{Regulation 12 requires that floors and traffic routes are constructed and maintained to enable safe movement; where these Regulations apply to a construction workplace (office, welfare facility, permanent access route), the duty to maintain surfaces and traffic routes in a suitable condition applies.}
  \bitem{MHSWR 1999}{Requires that risk assessments identify slips, trips and falls as specific hazards for each work activity and location; the assessment must consider the specific conditions under which work is carried out (weather, time of day, lighting, footwear) rather than applying a generic slip and trip statement.}
  \bitem{HSG150 Health and Safety in Construction}{Provides sector-specific guidance on housekeeping, site tidiness standards, management of temporary access routes, lighting requirements, and the organisation of materials storage to avoid obstruction of pedestrian and emergency access routes.}
\end{itemize}

%---------------------------- SECTION ----------------------------
\section{Conducting the assessment using the five-step method}

\subsection{Step 1 --- Identify hazards}
\paragraph{Slips, trips and fall hazards in construction include: wet, muddy, or oil-contaminated ground surfaces on access routes, hardstandings, and working platforms; uneven ground surfaces including pot-holes, ruts, material stockpiles, and excavation spoil in pedestrian routes; debris in circulation routes including wire ties, screed mesh, off-cuts, packaging, and construction waste; cables, hoses, and air lines crossing pedestrian routes; changes in level without marking or edge protection including steps between working areas, temporary ramps, and scaffold platform changes; inadequate or absent lighting on access routes, in covered areas, and during poor weather; unsuitable footwear worn by operatives, delivery personnel, and visitors; wet or frost-covered temporary access gangways and scaffold platforms; and the absence of defined, segregated pedestrian routes on sites where pedestrians and plant share the same ground.}

\subsection{Step 2 --- Decide who or what is affected}
\paragraph{Persons at risk from slips, trips and falls on construction sites include: all site operatives on access routes and in circulation areas; delivery drivers and haulage contractors accessing the site; specialist trade operatives who may not have received a full site induction covering the specific housekeeping and access standards; site visitors including clients, inspectors, and professional team members; and members of the public adjacent to hoarded sites where construction materials, debris, or liquids may affect the public footpath. Particular vulnerability applies to: workers with impaired balance or mobility; workers wearing unsuitable footwear; workers carrying loads that restrict their view of the ground ahead; and workers operating in low light or at night.}

\subsection{Step 3 --- Evaluate likelihood and consequence}
\paragraph{Slips, trips and falls on the same level carry a probability of occurrence that is High (3--4 on a 5×5 likelihood scale) on unmanaged construction sites due to the ubiquity of the hazard conditions. The consequence, while typically less severe than falls from height, is routinely assessed as Medium (fractures, sprains, lacerations) and can be High where the fall involves secondary contact with sharp or hard objects or where the operative is carrying a load that prevents them from arresting the fall. Inherent risk is typically rated as High (12--16 on a 5×5 matrix). Good housekeeping, defined pedestrian routes, and regular inspection can reduce residual risk to Low (4--6/25).}

\subsection{Step 4 --- Select and implement controls}
\paragraph{Controls for slips, trips and falls must begin with the physical environment: define and mark pedestrian routes on the Construction Phase Plan and on site; install proper temporary road surfaces (crushed stone, geotextile, or temporary road matting) on primary access routes to prevent mud and ruts; maintain defined separation between plant routes and pedestrian routes; establish and enforce a housekeeping standard including a specific responsibility matrix for each area of the site; install adequate temporary lighting in all covered, enclosed, or poorly-lit areas; establish cable management standards to prevent cables and hoses crossing pedestrian routes; enforce footwear standards including anti-slip, steel toe-capped safety boots as a minimum site requirement.}

\subsection{Step 5 --- Record and review}
\paragraph{Slips, trips and falls assessments must be reviewed as the site develops and the pedestrian routes, ground conditions, and activity patterns change. A site induction-phase assessment that addresses access routes through a cleared site will be inadequate once substructure work, drainage installation, and materials delivery have created the hazard landscape described above. Weekly site safety inspections must include a specific check of pedestrian routes, housekeeping standards, and lighting adequacy, with findings recorded and acted upon. All slip, trip and fall incidents must be investigated and their findings fed back into the assessment.}

\example{Five-Step Application}{A mechanical and electrical subcontractor is installing first-fix services in a partially-completed commercial building in January. The building shell is wind- and watertight but has no permanent lighting or finished floor surface. Step 1 identifies: wet concrete subfloor surfaces; cable runs from temporary power supplies crossing work routes; inadequate temporary lighting in service riser areas; uneven surface at ground-floor slab edge change. Step 2 identifies: M\&E operatives (8); apprentices (2); site manager conducting daily inspection. Step 3 rates inherent slip/trip risk as High (12/25). Step 4 implements: temporary LED string lighting installed throughout all work areas before M\&E operatives commence; cable management clips used to raise all temporary power cables to wall-level; wet surface areas marked with barriers and alternative routes defined; site rules enforced for safety boot requirement. Residual risk rated Low (6/25). Step 5: weekly inspection of cable management and lighting by site supervisor.}

%---------------------------- SECTION ----------------------------
\section{Applying the hierarchy of controls}

\paragraph{Slip, trip and fall controls follow the same hierarchy as all hazard controls, but the ubiquity of the hazard conditions in construction means that engineering controls (surface treatments, route design, lighting) must be planned before work commences rather than applied reactively after incidents.}

\begin{itemize}
  \bitem{Elimination}{Design the site layout to eliminate hazards in pedestrian routes: position material storage away from access routes; design delivery logistics to avoid materials being moved through populated work areas; specify ground-level prefabrication to reduce movement of loose materials through the site.}
  \bitem{Substitution}{Replace temporary earthen access tracks with crushed stone surfacing or temporary roadway matting; replace open-cable temporary power runs with wall-fixed or overhead distribution systems; replace smooth-surface temporary walkway boards with anti-slip mesh gangway panels.}
  \bitem{Engineering controls}{Install defined, drained, and surfaced pedestrian routes with proper edge marking; provide temporary anti-slip matting at building entrance and exit points; install adequate temporary lighting (minimum 100 lux on working areas, 50 lux on access routes) in all areas without natural light; install cable management systems to take all temporary cables to wall or overhead routes; provide handrails at all temporary steps and level changes.}
  \bitem{Administrative controls}{Publish and enforce site housekeeping standards including responsibility allocation for each area; conduct weekly site safety inspections with specific slip, trip and fall checklist; enforce footwear standards at site induction and maintain a register of footwear spot-checks; include slip, trip and fall hazards in trade-specific toolbox talks relevant to each phase; establish and publicise an on-site waste disposal system to prevent debris accumulation.}
  \bitem{PPE}{Anti-slip, steel toe-capped safety boots conforming to EN ISO 20345 as a minimum site requirement; high-visibility vests to ensure workers are visible during plant movement near pedestrian routes; knee pads for operatives working at low level to reduce injury severity in the event of a fall.}
\end{itemize}

%---------------------------- SECTION ----------------------------
\section{Cadence of assessment and review triggers}

\paragraph{Slip, trip and fall conditions on a construction site change continuously as work progresses. The assessment cannot be treated as a one-off exercise; it must be incorporated into the regular site safety management routine.}

\begin{itemize}
  \bitem{Weekly site safety inspection}{Every site safety inspection must include a specific walk of pedestrian routes, access points, and working areas with a focus on housekeeping, lighting, cable management, and surface condition; findings must be recorded and deficiencies allocated for remedy with a defined deadline.}
  \bitem{After any weather event}{Rain, frost, or snow significantly increases slip risk on all outdoor access routes; following any significant weather event, access routes must be inspected before work commences and any areas of standing water, ice, or mud that obstruct routes must be treated before operatives move through them.}
  \bitem{At the start of each new work phase}{When a new trade commences work, the pedestrian routes through their work area must be reviewed; new trades bring new materials, new cable runs, and new working patterns that may create hazards not present in the previous phase.}
  \bitem{Following any slip, trip or fall incident}{All slip, trip and fall incidents, however minor, must be investigated and the findings used to update the assessment; a trend of minor incidents in a particular area indicates a systemic control failure.}
  \bitem{When site layout or access routes change}{Any significant change in site layout, welfare facility position, material storage areas, or pedestrian access routes must trigger a review of the slip, trip and fall assessment for the affected areas.}
\end{itemize}

%---------------------------- SECTION ----------------------------
\section{Common failure modes}

\begin{itemize}
  \bitem{Housekeeping responsibility not allocated}{``Good housekeeping'' is specified as a control in every RAMS document on every construction site, but responsibility for maintaining it in specific areas is rarely allocated in writing. The result is that no individual takes ownership and the standard is never enforced.}
  \bitem{Pedestrian routes not defined or not surfaced}{On many construction sites, there are no defined pedestrian routes through the working areas; operatives pick their own routes through the most hazardous conditions rather than following a planned and controlled path.}
  \bitem{Temporary lighting not provided or not maintained}{Temporary lighting lamps fail, are knocked down, or are removed and not replaced. Areas that were adequately lit at the start of a work phase become progressively darker as the phase extends and bulb failures accumulate without replacement.}
  \bitem{Cable management absent}{Temporary power distribution by means of cables draped across floors and gangways is ubiquitous on construction sites and is one of the most common trip hazards identified in site inspections and incident investigations.}
  \bitem{Footwear standards not enforced at gate}{Site induction footwear rules are stated but not enforced at the site entrance; delivery drivers, sub-subcontractors, and visitors routinely access the site in trainers or unsuitable footwear, and no mechanism exists to enforce the standard.}
  \bitem{Incidents treated as minor and not investigated}{Slips and trips that result in minor injuries are frequently treated as trivial events that do not require investigation. The pattern of minor incidents is the leading indicator of the serious injury that is statistically inevitable if the hazard conditions remain unchanged.}
\end{itemize}

\example{Failure Pattern}{An electrical contractor is running conduit on a multi-storey residential scheme during the first-fix phase. Temporary power is supplied via four long extension leads draped from the supply point, across the corridor, and into each room being worked on. The leads are not fixed to the walls or overhead and lie across the floor of the central corridor used by all trades. Over the course of the first-fix phase, three tripping incidents are recorded, each resulting in minor injuries treated on site. None is reported to the site manager; none is investigated. In week seven, a plasterer trips on a cable while carrying a stilts frame, falls, and sustains a fractured elbow. The investigation reveals seven trip incidents involving the same cables over eight weeks. The cables had been identified in the electrical contractor's RAMS as requiring management; no management had been implemented.}

%---------------------------- CHAPTER SUMMARY ----------------------------
\section{Chapter Summary}

\begin{table}[H]
\centering
\begin{tabularx}{\textwidth}{>{\raggedright\arraybackslash}p{3.2cm}X}
\toprule
\textbf{Assessment Element} & \textbf{Operational Requirement} \\
\midrule
Legal/Professional Basis & MHSWR 1999; CDM 2015; Workplace (Health, Safety and Welfare) Regulations 1992. Principal contractor must plan pedestrian routes and housekeeping standards in the Construction Phase Plan. \\
Method & Apply the full five-step process and document inherent/residual risk. \\
Controls & Defined and surfaced pedestrian routes; housekeeping standards with allocated responsibility; adequate temporary lighting; cable management systems; enforced footwear standards. \\
Cadence & Weekly site safety inspection; review after adverse weather, phase changes, and any incident. \\
Assurance & Use audit, incident learning, and governance reporting to verify effectiveness. \\
\bottomrule
\end{tabularx}
\end{table}

\begin{itemize}
  \bitem{Key takeaway 1}{Slips, trips and falls are the most frequent cause of major injury on construction sites; they are preventable through site organisation, not through reminders to watch where you are walking.}
  \bitem{Key takeaway 2}{Housekeeping as a control is only effective when responsibility is allocated in writing to a named individual or role for each area of the site; generic ``good housekeeping'' in a RAMS has never prevented a single injury.}
  \bitem{Key takeaway 3}{Cable management is a primary engineering control for trip hazards and must be designed into temporary power distribution plans from the outset rather than treated as a housekeeping matter.}
  \bitem{Key takeaway 4}{Minor slip, trip and fall incidents are leading indicators of future serious injuries; each must be investigated and the findings used to update the assessment and improve controls.}
\end{itemize}

\barquote{In Chapter~\ref{chap:Plant, Machinery and Vehicle Risk Assessment}, we turn from this assessment to the next domain in the Prudentia framework.}
